\chapter{后端引擎}\label{chap:backend}

本章描述系统的后端计算引擎，由四个紧密耦合的组件组成：异构图建模将原始档案重写为带类型边的图、异构图变换器在该图上学习节点表征并对候选婚配对打分、子图链路预测从候选对周围的封闭子图抽取局部结构指纹、多智能体协商协议把上述证据"重演"为可审计的推理过程。四者按"图—编码器—子图—协商"次序衔接，向第~\ref{chap:vis}~章的可视分析系统输出排序好的候选对与全套证据轨迹。

\section{异构图建模}\label{sec:backend-graph}

在切入异构图形式化之前，需先厘清本研究面对的两类档案对象。一类是带标注训练源 A，对应基于辽宁地区清代户籍档案整理的中国多代人口面板（China Multi-Generational Panel Dataset，CMGPD-LN）\citep{campbell2011data,zhao2001chinese}，覆盖 1749—1909 年八个普查年共约 1{,}513{,}357 条人年观测，对应约 266{,}091 个唯一个体与约 85{,}538 条可识别婚姻边；其特殊之处在于地方衙门按男性当家人编册时已把入籍女性以"妻"位与"母"位同步登记，从而保留了父系、母系、共居、村落与八旗等多层级结构信息，并提供了可作训练监督的婚姻边真值。另一类是目标域 B，对应历史学家在田野调查中收集到的真实家谱集合：B 多为父系单线编纂、女性条目体例参差、跨家族的婚姻关系几乎从不显式记录，因而在图论意义上呈现为一片彼此孤立的父系树森林，缺少婚姻边的完整真值。

A 与 B 是同一历史现象的两个侧面：A 提供了少量被户籍登记意外补全的样本，B 则承载了下游应用真正待补的对象。直接在 B 上做监督学习不可行，因为 B 本身没有真值；直接把 A 上学到的模型迁移到 B 又面临两条理想化前提。其一为封闭性假设，即同一份 B 内的男性候选与同一份 B 内的女性候选穷尽对应，不存在 B 内男性与 B 外女性的婚姻；其二为特征对齐假设，即 B 与 A 在人物特征空间上一致或可线性对齐。两条假设在真实田野家谱中都只能近似成立，但承认这一近似才能把 A→B 的开放问题简化为可在受控实验中度量的统计学习问题。

基于此，本研究在 A 上构造一个简化代理：保留 A 的全部节点与多层级关系，但人为遮蔽与目标域 B 同形的母系登记，再在该简化代理上训练评分模型，从而在已知真值的条件下度量结构信号的贡献。该简化代理与目标域 B 之间存在系统性差距，行文中将一以贯之地以"简化代理"称呼经处理的 A，避免与"等价于 B"的更强主张混淆。

将上述档案形式化为异构图 \(G=(V,E,\mathcal{T},\mathcal{R})\)：\(\mathcal{T}=\{\text{个体},\text{家户},\text{村落},\text{八旗}\}\)；\(\mathcal{R}\) 共九类，按语义归为三组：父系亲属（\(r_{fs}\) 父→子、\(r_{fd}\) 父→女）、母系亲属（\(r_{ms}\) 母→子、\(r_{md}\) 母→女）、共父同胞 \(r_{sib}\)、共居与隶属（\(r_{hh}\) 人→家户、\(r_{hc}\) 家户→村落、\(r_{cb}\) 人→八旗）以及带时间戳的婚姻边 \(r_{hw}\)（夫指向妻）。亲属边的时间戳取子女出生年，家户隶属边取观测年，村落与八旗等空间隶属边视为静态边、时间戳取 0。完整登记见表~\ref{tab:edge-types}。

\begin{table}[h]
\centering
\caption{异构图的九类边类型登记。时间戳列中"是"代表该类边携带逐边时间戳，"否（静态）"代表全部边时间戳取 0 而在任意年份子图中均保留。}
\label{tab:edge-types}
\begin{tabular}{llll}
\hline
名称 & 起止节点类型 & 语义 & 时间戳 \\
\hline
$r_{hw}$ & 个体 $\to$ 个体 & 婚姻边，夫指向妻 & 是（成婚年） \\
$r_{fs}$ & 个体 $\to$ 个体 & 父亲指向儿子 & 是（子女出生年） \\
$r_{fd}$ & 个体 $\to$ 个体 & 父亲指向女儿 & 是（子女出生年） \\
$r_{ms}$ & 个体 $\to$ 个体 & 母亲指向儿子 & 是（子女出生年） \\
$r_{md}$ & 个体 $\to$ 个体 & 母亲指向女儿 & 是（子女出生年） \\
$r_{sib}$ & 个体 $\to$ 个体 & 共父同胞，双向连接 & 是（较晚出生年） \\
$r_{hh}$ & 个体 $\to$ 家户 & 个体当年所在家户 & 是（观测年） \\
$r_{hc}$ & 家户 $\to$ 村落 & 家户所在村落 & 否（静态） \\
$r_{cb}$ & 个体 $\to$ 八旗 & 个体所属八旗 & 否（静态） \\
\hline
\end{tabular}
\end{table}

异构图 \(G\) 是面板期内全部观测的叠加，若直接送入图神经网络做链路预测，会让早年队列的婚姻边经过晚年才出现的边间接获得未来信息。为消除此类泄漏，本研究为每个目标年 \(t\) 构造时间因果子图：
\begin{equation}
G_t = \bigl(V,\ \{e\in E : \mathit{edge\_time}(e) < t\},\ \mathcal{T},\ \mathcal{R}\bigr).
\label{eq:subgraph}
\end{equation}
即在不改变节点与类型集合的前提下只保留发生时间严格早于 \(t\) 的有向边。静态边时间戳为 0，对任意 \(t\ge 1\) 由 \(0<t\) 恒成立可知它们在所有时间因果子图中始终保留。训练时还在消息传递图中额外移除当前队列年的目标对以及验证集、测试集的全部正样本对。

为在 A 上构造与目标域 B 同形的简化代理，需要把 A 中那部分"超出 B 通常记载范围"的母系登记从图中遮蔽掉。此即父系镜像消融：对边类型 \(r\in\{r_{ms},r_{md}\}\) 保留键于元数据快照中、将其邻接索引清空为零长度。之所以坚持"消融而不删除"，是因为异构图变换器按元数据快照为每条关系注册独立的线性投影矩阵\citep{hu2020hgt}；若运行期移除某关系键，与之绑定的投影权重便会与图结构错位、使训练好的快照失效。保留键、置空索引只让该关系传递空消息而不参与聚合，既等价于物理删除信号，又规避元数据漂移。

须强调父系镜像消融与目标域 B 不构成等价镜像，而是简化代理。真实田野家谱的女子入谱体例随时代与族属差异显著：部分时期女性仅以"妻"位入夫家而不进入母家谱系，部分时期女儿、妻、妾的特征记录极稀疏甚至只剩占位姓氏\citep{yan2009nuzi}；而 A 的母系边在被消融后所留下的空键并不能完整复刻这种参差。同样地，封闭性假设与特征对齐假设在简化代理中都被默认满足，而在 B 上则只能近似成立。该简化代理相对真实场景的差距将在 \S\ref{sec:conclusion-limit}~系统讨论；本节后续讨论限定在简化代理之上，以便公式与指标都有明确的图论与统计学意义。

\section{异构图变换器编码与婚姻评分}\label{sec:backend-hgt}

人物节点是评分任务唯一直接参与的节点类型。为其设计五条平行特征流：性别向量（保留 1=女、2=男的原始编码）、关系嵌入、出生年的列内 \(z\) 分数、官职品阶俸禄三项 0/1 指示位，以及宏观协变量（同区域粮价 \(z\) 分数与逐年差分、清代纪元号、灾害指示、队列年 \(z\) 分数）。五条流拼接后由一层线性投影统一映射到节点向量。宏观协变量不写入静态特征而由队列年驱动逐次广播：时间因果子图构造时按目标年读取对应行挂到当次前向使用的人物特征上，使任意快照年看到的宏观信号都严格对齐当前打分的队列年。家户、村落、八旗三类节点分别以户口规模、村庄人口、旗别独热等少量字段经同构投影映射到同一隐空间，便于跨类型注意力。

设节点 \(v\) 在第 \(\ell\) 层的隐向量为 \(h_v^{(\ell)}\)，节点类型函数为 \(\tau:V\to\mathcal{T}\)，关系类型集合为 \(\mathcal{R}\)。异构图变换器在每条边类型 \(r\in\mathcal{R}\) 上以多头注意力机制\citep{vaswani2017attention}聚合邻居信息，并通过类型化的查询、键、消息投影区分不同节点类型与不同关系\citep{hu2020hgt}。注意力分数定义为
\begin{equation}\label{eq:hgt-att}
  \alpha_{u,v}^{r} \;=\; \mathrm{softmax}_u\!\left(
    \frac{\bigl(W^{q}_{\tau(v)}\,h_v^{(\ell)}\bigr)^{\!\top}\,W^{r}_{\mathrm{att}}\,\bigl(W^{k}_{\tau(u)}\,h_u^{(\ell)}\bigr)}
         {\sqrt{d/H}}
    \cdot \mu_{\tau(u),\,r,\,\tau(v)}\right),
\end{equation}
节点更新规则为
\begin{equation}\label{eq:hgt-upd}
  h_v^{(\ell+1)} \;=\; \mathrm{LN}\!\left(
    h_v^{(\ell)}\;+\;\mathrm{Dropout}_p\!\Bigl(
      \sum_{r\in\mathcal{R}}\sum_{u\in\mathcal{N}_r(v)}
        \alpha_{u,v}^{r}\,W^{r}_{\mathrm{msg}}\,W^{v}_{\tau(u)}\,h_u^{(\ell)}\Bigr)
    \right).
\end{equation}
其中 \(W^{q}_{\tau(v)}\)、\(W^{k}_{\tau(u)}\)、\(W^{v}_{\tau(u)}\) 按节点类型注册；\(W^{r}_{\mathrm{att}}\) 与 \(W^{r}_{\mathrm{msg}}\) 按关系注册；\(\mu\) 为可学习的关系先验张量，给"来源类型—关系—目标类型"三元组提供全局乘性偏置。残差项与外层层归一化共同稳定深层堆叠下的梯度；丢弃以概率 \(p\) 在训练阶段对聚合后的消息随机置零以缓解过拟合\citep{kipf2016gcn}。

编码器输出每个人物节点隐表征 \(h\) 后，由一个轻量二元交互多层感知机给候选对 \((m,w)\) 打分，沿用四种线性无法捕捉的二阶交互形式\citep{vaswani2017attention}：
\begin{equation}\label{eq:scorer}
\begin{aligned}
  \mathbf{p}_{m,w} &= [h_m \,\|\, h_w \,\|\, |h_m-h_w| \,\|\, h_m\odot h_w]\in\mathbb{R}^{4d}, \\
  \psi(h_m,h_w) &= W_2\,\mathrm{Dropout}_p\bigl(\mathrm{GELU}(W_1\,\mathbf{p}_{m,w})\bigr)\;\in\;\mathbb{R}.
\end{aligned}
\end{equation}
原始两向量保留全局位置，绝对差编码两者在嵌入空间的对称差异，哈达玛积编码逐维一致性；四种交互拼接后由 \(W_1\) 投到隐空间，再由 \(W_2\) 输出标量 logit。

训练循环按队列年为粒度组织：每步随机抽取若干队列年合并为一次优化器步，对每个队列年先按式~\eqref{eq:subgraph}~构造时间因果子图，再扣除当前队列年的目标对以及验证集、测试集的全部正样本对以避免标签泄漏。负样本由"队列内负采样"给出——同一队列年内通过替换女性候选生成反例；跨年负采样会让模型只需识别队列年差异即可降低损失，使排序指标显著虚高，但对"同一年内的若干合理候选间排序"无信号增益。每一批次的损失为
\begin{equation}\label{eq:loss}
\begin{aligned}
  \mathcal{L} &= -\frac{1}{|\mathcal{B}|}\sum_{(m,w,y)\in\mathcal{B}}\ell(m,w,y), \\
  \ell(m,w,y) &= y\log\sigma(\psi(h_m,h_w)) + (1-y)\log\bigl(1-\sigma(\psi(h_m,h_w))\bigr).
\end{aligned}
\end{equation}
\(\sigma\) 为 sigmoid，\(\mathcal{B}\) 为该批次中所有正负样本三元组。优化器为 AdamW 配合余弦退火学习率，整体基于 PyTorch\citep{paszke2019pytorch} 与 PyG\citep{fey2019pyg} 实现。

数据按婚姻发生年切分：不晚于 1855 进入训练集、到 1879 进入验证集、到 1909 进入测试集。验证阶段对每个队列年做一次匈牙利最优一对一指派并统计 recall@1，该指标比 ROC-AUC 更贴近档案补全的实际诉求。验证最优的训练快照分别为消融条件与未消融条件保存，两份快照不可互换：消融图缺少两条母系关系，因此其训练得到的母系关系投影从未被更新，载入未消融图则相当于在随机初始化关系上推断，结果失去意义。最终在 1880—1909 留出测试集（共 131{,}140 个候选对）上，编码器与评分多层感知机共同给出 \(\mathrm{Hit@10}=0.624\)、\(\mathrm{MRR}=0.428\)；母系边消融对照下，未消融上界相对消融模型在按队列年聚合的匈牙利 recall@1 上平均高出约 \(62.6\%\)，说明母系边在父系碎裂图上对婚姻边的可识别度仍有显著贡献。

\section{子图链路预测}\label{sec:backend-seal}

整图层面的婚姻评分给出可比较的全局信号，但单一标量难以提供"为何此对会被推荐"的局部依据。为此本节另起一支独立分支，把每个候选对周围的局部拓扑显式提取出来并以可命名的微结构图样形式向下游协商提供结构指纹。给定候选对 \((m,w,\tau)\)，子图链路预测（Subgraph Embedding And Labeling，SEAL）\citep{zhang2018seal} 把"\(m\) 与 \(w\) 在 \(\tau\) 年是否结为婚配"重新表述为以二元组为中心的封闭子图二分类：从因果可见图（其上母系边按 \S\ref{sec:backend-graph}~的镜像消融策略置空但保留元数据键）出发，沿无向投影抽取以 \(m,w\) 为双中心的 \(h\) 跳封闭子图，再依次施加六条反泄漏屏蔽以删除可能直接揭示婚姻标签的边，得到掩码子图。子图判据被约束在直径不超过 \(h\) 跳的局部邻域内，便于映射回宗族、户籍、八旗等可命名的史学概念；该路得分将在 \S\ref{sec:backend-mas}~的协商协议中与编码器全局得分作为独立证据进入合议。

封闭子图剥离了原图的全局坐标，模型必须从结构本身重新获得"哪些节点扮演什么角色"。为此引入双半径节点标注（Double-Radius Node Labeling，DRNL）：对子图内任一节点 \(v\)，记 \(d(v,m)\) 与 \(d(v,w)\) 为 \(v\) 到两端点的无向最短路径长度（\(\infty\) 表示不可达），令 \(d=d(v,m)+d(v,w)\)，则
\begin{equation}\label{eq:drnl}
\ell(v)=
\begin{cases}
0, & d(v,m)=\infty\ \text{或}\ d(v,w)=\infty,\\[3pt]
1, & v=m\ \text{或}\ v=w,\\[3pt]
1+\min\!\big(d(v,m),d(v,w)\big)+\dfrac{d}{2}\!\left(\dfrac{d}{2}+d\bmod 2-1\right), & \text{其他}.
\end{cases}
\end{equation}
标号对距离对严格对称，交换男女两端不改变结果；同距离对的节点共享同一标号，从而把节点离散化为有限多个"距离 \(\times\) 拓扑组合"等价类。在掩码子图上枚举从 \(m\) 到 \(w\) 的简单路径并把途径边按可见关系拼接为序列，所得序列称为一个微结构图样。在 CMGPD-LN 训练集 80 个查询事件上筛选出 13 类史学含义明确、识别频率较高的图样，按"父系结构""户与姻亲网络""八旗结构"三组归类，完整目录列于表~\ref{tab:motif-catalog}。

\begin{table}[ht]
  \centering
  \caption{十三类历史微结构图样目录。覆盖率为掩码子图上 80 个训练事件中至少存在一条对应路径的比例。}\label{tab:motif-catalog}
  \begin{tabular}{c l l c}
    \toprule
    编号 & 关系序列 & 史学含义 & 覆盖率（\%）\\
    \midrule
    \multicolumn{4}{l}{\emph{父系结构}}\\
    M01 & $(r_\mathrm{sib})$ & 同父登记的兄妹 & 86.7 \\
    M02 & $(r_\mathrm{fs},r_\mathrm{fd})$ & 共同父亲（最关键信号） & 93.3 \\
    M03 & $(r_\mathrm{sib},r_\mathrm{sib})$ & 堂表兄妹 & 80.0 \\
    M04 & $(r_\mathrm{sib},r_\mathrm{sib},r_\mathrm{sib})$ & 大宗族网络 & 80.0 \\
    M05 & $(r_\mathrm{fs},r_\mathrm{fs},r_\mathrm{sib})$ & 父系祖辈兼兄妹（辈分错位） & 80.0 \\
    M06 & $(r_\mathrm{sib},r_\mathrm{fs},r_\mathrm{fd})$ & 共父展开变体 & 80.0 \\
    M07 & $(r_\mathrm{fs},r_\mathrm{fd},r_\mathrm{sib})$ & 共父再延一步兄妹 & 73.3 \\
    M08 & $(r_\mathrm{sib},r_\mathrm{fd},r_\mathrm{fd})$ & 代际错位的同宗联姻链 & 73.3 \\
    \multicolumn{4}{l}{\emph{户与姻亲网络}}\\
    M09 & $(r_\mathrm{hh},r_\mathrm{hh},r_\mathrm{sib})$ & 同户成员之兄妹（包户亲属） & 50.0 \\
    M10 & $(r_\mathrm{hh},r_\mathrm{hh},r_\mathrm{fd})$ & 同户成员之女（主家女儿） & 66.7 \\
    M11 & $(r_\mathrm{sib},r_\mathrm{hw},r_\mathrm{sib})$ & 前代联姻链（重复联姻） & 33.3 \\
    \multicolumn{4}{l}{\emph{八旗结构}}\\
    M12 & $(r_\mathrm{fs},r_\mathrm{cb},r_\mathrm{cb})$ & 八旗内同父系（旗内婚父系视角） & 83.8 \\
    M13 & $(r_\mathrm{cb},r_\mathrm{cb})$ & 同旗（旗内婚） & 85.0 \\
    \bottomrule
  \end{tabular}
\end{table}

掩码子图被送入一个四层关系图卷积网络（Relational GCN，R-GCN）评分器。节点输入特征由节点类型独热、双半径标号独热与若干人口学数值特征拼接而得；卷积更新满足
\begin{equation}\label{eq:rgcn-update}
h_v^{(\ell+1)}=\sigma\!\left(\sum_{r\in\mathcal{R}_\mathrm{vis}}\sum_{u\in\mathcal{N}_r(v)}\frac{1}{c_{v,r}}\,W_r^{(\ell)}\,h_u^{(\ell)}+W_0^{(\ell)}\,h_v^{(\ell)}\right),\quad W_r^{(\ell)}=\sum_{b=1}^{B}a_{rb}^{(\ell)}V_b^{(\ell)}.
\end{equation}
其中 \(\mathcal{N}_r(v)\) 为 \(v\) 沿关系 \(r\) 的邻居集合，\(c_{v,r}\) 为标准化常数，\(\sigma\) 取 ReLU。关系特定矩阵经基分解由 \(B=4\) 个共享基矩阵线性组合而得；由于基数严格小于可见关系数，不同关系被强制共享部分参数，缓解关系数偏多带来的过拟合。读出阶段采用 SortPooling 与两层小型卷积加多层感知机映射为标量得分，整体可学习参数 97{,}889。

训练采用列表式 softmax 交叉熵：对每个真实婚姻事件，从同期合法候选池中按"10 个同社区 + 10 个全局"策略采集 20 个负样本，最小化对真实婚配的负对数似然；同社区难负样本迫使模型不仅依赖社区粗粒度信号。在 80 个查询的测试集上，模型在 100 候选随机配对池中给出 Hit@1 \(=0.975\)、MRR \(=0.984\)，相较随机基线提升约 97 倍。更具诊断意义的是按男性父系链长度分组的消融：当登记父亲缺失时，Hit@1 由父亲已知组的 \(1.000\) 急剧下滑至 \(0.429\)，而仅父亲、祖父及更长链长组合并的群体仍稳定达到 \(1.000\)。这一差异说明"局部结构指纹"的主导成分仍是父系结构信号——父系链断裂后子图内可被关系图卷积网络利用的判据急剧减少，与 \S\ref{sec:backend-hgt}~母系边消融的结论方向一致。

\section{多智能体协商协议}\label{sec:backend-mas}

前两节给出的是以分数呈现的候选婚配对，单一分数难以替代历史学家的权衡式判断。本节将结构化证据"重演"为协商过程，借鉴生成式智能体\citep{park2023generative} 与思维链\citep{wei2022cot} 并在多轮辩论\citep{liang2023llm} 的交互结构上特化。鉴于已有研究\citep{turpin2023faithfulness} 指出大语言模型存在事后合理化倾向，本协议限制智能体自由度：所有论据须挂载外部证据，每轮陈述与评分完整记录以备审计。

对于队列年中每个候选对 \((m,w)\)，协议实例化夫方与妻方两个智能体，由同一模板生成但承担互补的论辩立场。模板由六个中文字段构成：人物档案（姓名、生年、性别、官职、八旗归属与户位编码，缺失时渲染为"档案未载"）；异构图变换器预分（由 \S\ref{sec:backend-hgt}~的评分 logit 映射到 \([0,10]\) 区间并附置信带）；生命事件块（从生命事件主表读取该年的第一事件代码与第二事件代码并经事件代码翻译词典渲染为可读短语，在册而无事件时与"档案未载"严格区分）；收入桶（按百分位划分为低于 33 分位、33—66 分位、高于 66 分位三档）；子图模式标签（列出 \S\ref{sec:backend-seal}~中命中的微结构图样，每一命中以"可读名称（双半径标号摘要）"呈现）；提示（承载历史学家通过控制台累计注入的消息，按地址 tag 路由）。

整套协商组织为六轮：印象生成、深度询问、反驳、对齐、复述与终评，借鉴双方协商中"先建立印象、再追问差异、继而互相挑战、最后达成或保留分歧"的路径，并在轮间保留历史学家手动推进的钩子。算法~\ref{alg:negotiation} 给出伪代码。

\begin{algorithm}[t]
\caption{六轮多智能体协商：候选集合 \(\mathcal{C}\) 在队列年 \(t\) 的协商流程。}\label{alg:negotiation}
\begin{algorithmic}[1]
\Require 候选集合 \(\mathcal{C}\)；提示队列 \(Q\)；推进信号 \(\mathcal{E}\)；最终得分超参 \(\lambda=0.30\)
\Ensure 按降序排列的候选与最终得分对 \(\{(c, \mathrm{score}_i)\}\)
\State \textbf{第 1 轮（印象生成）：}对 \(c \in \mathcal{C}\) 通过并行调度原语并行实例化夫方与妻方智能体，输出履历段落与初评分 \(s_i^{(1)}, t_i^{(1)} \in [0,10]\)
\For{第 \(r\) 轮 \(r \in \{2,3,4,5\}\)（依次对应深度询问、反驳、对齐、复述）}
  \State 从 \(Q\) 取出本轮提示，按地址 tag \texttt{@everyone}、\texttt{@target}、\texttt{@c-XX} 路由
  \State 将对方上一轮陈述与新提示注入两侧上下文
  \State 两侧智能体输出更新陈述与评分 \(s_i^{(r)}, t_i^{(r)}\)
  \State 等待轮间推进信号 \(\mathcal{E}\)，由历史学家在界面按下"下一轮"按钮触发
\EndFor
\State \textbf{第 6 轮（终评）：}对每一 \(c \in \mathcal{C}\)，按式~\eqref{eq:final-score} 计算 \(\mathrm{score}_i\)，并按降序输出到前端
\State \Return \(\{(c, \mathrm{score}_i)\}_{c \in \mathcal{C}}\)
\end{algorithmic}
\end{algorithm}

第 1 轮通过并行调度原语并行实例化两端智能体；第 2 至 5 轮在每轮起始先排空提示队列，把命中提示与对方上一轮陈述合并到上下文再请求新陈述与评分；轮间推进信号使协商不会一气呵成连跑到底，历史学家可在每轮末端复核再决定是否推进。第 6 轮仅按下式合成最终得分并输出排序，将语义生成与分数合成显式分离。记夫方第 5 轮终态评分为 \(t_i\in[0,10]\)、妻方对位评分为 \(s_i\in[0,10]\)，候选对 \(i\) 的最终得分定义为
\begin{equation}\label{eq:final-score}
  \mathrm{score}_i = \frac{s_i + t_i}{2} - \lambda\,|s_i - t_i|, \quad \lambda = 0.30.
\end{equation}
前一项奖励双方一致高分；后一项以绝对差惩罚不对称同意。例如 \((9,9)\) 仍得 9、而 \((9,5)\) 降至 5.8 低于均值 7。\(\lambda=0.30\) 由形成性研究与专家共同确定：0.5 以上则合理候选被压至 5 分以下失去区分度，0.1 以下则不对称同意几乎不被惩罚，与"双边匹配应彼此说服"的史学直觉相悖。

历史学家在轮间可通过控制台注入提示。每条提示由一个地址 tag 与一个动词组成，前者决定作用域、后者决定智能体如何调整评分倾向，语义见表~\ref{tab:hint-routing}。

\begin{table}[h]
\centering
\caption{提示路由与动词语义。}\label{tab:hint-routing}
\footnotesize
\begin{tabular}{@{}lll@{}}
\toprule
地址 tag / 动词 & 作用域 & 对智能体推理的影响 \\
\midrule
\texttt{@everyone} & 当前队列年所有活跃智能体 & 提示广播到所有候选对的两侧角色 \\
\texttt{@target} & 当前选中候选对的两侧智能体 & 提示仅路由到选中对的夫方与妻方 \\
\texttt{@c-XX} & 编号为 XX 的候选对的两侧智能体 & 提示仅路由到该具体候选对 \\
\midrule
\texttt{boost} & 与上述任一作用域配合 & 指示智能体在下一轮提高对该候选的评分 \\
\texttt{penalise} & 与上述任一作用域配合 & 指示智能体在下一轮降低对该候选的评分 \\
\texttt{eliminate} & 与上述任一作用域配合 & 指示智能体将评分置 0 并在后续轮次不再为其辩护 \\
\texttt{accept} & 与上述任一作用域配合 & 指示智能体将评分置 10 并在后续轮次强力支持 \\
\bottomrule
\end{tabular}
\end{table}

智能体收到提示后须在下一轮陈述中显式引用，既保证人为干预可被审计，也避免提示在长上下文中被悄然遗忘。协商所用语言模型为千问 Plus 模型（具体版本），通过其 OpenAI 兼容端点访问、凭据由大模型服务凭证注入。为使协议在缺乏凭据的开发与教学场景下仍可演示，本研究设计本地占位智能体兜底：凭据为空或调用失败时切换到占位智能体，每轮返回固定占位陈述与一个由初评分加小幅扰动得到的分数，使前端轮次推进、提示路由、最终评分与可视化均能照常运转；占位路径在前端以"模拟模式"显式标注，避免被误读为真实推理结果。
